\subsection{Irreducible polynomials}

DEFINITION 2. --- Let K be a commutative field. We say that \(f \in \mathbf{K}[\mathbf{X}]\) is irreducible if \(\deg f \geqslant 1\) and \(f\) is not divisible by any polynomial \(g\) such that \(0 < \deg g < \deg f\) .

It comes to the same to say that \(\deg f \geqslant 1\) and the only divisors of f in K{[}X{]} are the scalars \(\neq 0\) and the products of f by scalars \(\neq 0\) . Since the relation \((f) \subset (g)\) means that g divides f, we see that the irreducible polynomials of K{[}X{]} may also be defined as the polynomials f such that the ideal \((f)\) is maximal (I, p.~104).

Let \(f, g \in \mathbf{K}[\mathbf{X}]\) . If \(f\) is irreducible, it is clear that either \(f\) and \(g\) are relatively prime or that \(f\) divides \(g\) . If \(f\) and \(g\) are irreducible, either \(f\) and \(g\) are relatively prime or each is a product of the other by a scalar \(\neq 0\) . In particular, two distinct irreducible monic polynomials are relatively prime.

PROPOSITION 13. --- Let \(\mathcal{I}\) be the set of irreducible monic polynomials in K{[}X{]}. Let f be a non-zero element of K{[}X{]} and α its leading coefficient; then there exists precisely one family of positive integers \((v_{p})_{p\in \mathcal{I}}\) with finite support, such that we have a decomposition

\[
f = \alpha \prod_ {p \in \mathcal {I}} p ^ {v _ {p}}.\tag{18}
\]

It suffices to prove the proposition when \(f\) is monic, that is when \(\alpha = 1\) . We shall argue by induction on the degree \(n\) of \(f\) , the case \(n = 0\) being trivial. Suppose then that \(n \geqslant 1\) and that the proposition has been established for all polynomials of degree \(< n\) .

Let E be the set of monic polynomials \(\neq1\) which divide f; we have \(f\in E\) hence E is not empty and there exists in E a polynomial g of least degree. It is clear that g is irreducible and there exists a monic polynomial h of degree \textless n such that f=gh; by the induction hypothesis h is the product of a finite family of irreducible monic polynomials, hence f has the same property. This proves the existence of the decomposition (18).

Let us now prove the uniqueness of the decomposition (18). Let \((w_{p})_{p\in \mathcal{I}}\) be a family of positive integers with finite support, such that \(f = \prod_{p\in \mathcal{I}}p^{w_p}\) . Since \(f\) is of

degree \(n \geqslant 1\) , there exists \(p \in \mathcal{I}\) such that \(w_{p} > 0\) ; if we had \(v_{p} = 0\) , f would be the product of a family of elements of \(\mathcal{I}\) distinct from p, hence it would be prime to p (IV, p.~13, Cor. 6), contrary to the fact that p divides f.~By the induction hypothesis the polynomial f/p admits a unique decomposition of type (18); hence we conclude the equality \(w_{q} = v_{q}\) for every \(q \in \mathcal{I}\) .

Let \(f\) be a non-zero polynomial in \(\mathbf{K}[\mathbf{X}]\) . We shall say that \(f\) has no multiple factors if the exponents \(v_{p}\) in the decomposition (18) are all \(\leqslant 1\) ; it comes to the same to say that \(f\) is the product of a finite sequence of pairwise distinct irreducible polynomials, or also that \(f\) is not divisible by the square of any non-constant polynomial of \(\mathbf{K}[\mathbf{X}]\) .

